\documentclass{article}
\usepackage[left=2cm, right=2cm, top=2cm, bottom=2cm]{geometry}
\usepackage{graphicx}
\usepackage{amsmath}
\usepackage{amssymb}
\usepackage{amsthm}
\usepackage{fancyhdr}
\usepackage{verbatim}
\usepackage{listings}
\usepackage{xcolor}
\usepackage{pdfpages}
\usepackage{cancel}
\usepackage{physics}

\lstset{
    language=C++,
    basicstyle=\ttfamily\footnotesize,
    keywordstyle=\color{blue},
    commentstyle=\color{green},
    stringstyle=\color{red},
    numbers=left,
    numberstyle=\tiny\color{gray},
    frame=single,
    breaklines=true,
    captionpos=b,
}


\title{MTH 553: Lecture \# 19}
\author{Cliff Sun}

\newtheorem{theorem}{Theorem}[section]
\newtheorem{lemma}[theorem]{Lemma}
\newtheorem{definition}[theorem]{Definition}
\newtheorem{conjecture}[theorem]{Conjecture}
\newtheorem{proposition}[theorem]{Proposition}
\newtheorem{corollary}[theorem]{Corollary}
\newtheorem{one minute paper}[theorem]{One Minute Paper}

\pagestyle{fancy}
\lhead{\textbf{\thepage}\ \ \nouppercase{\rightmark}}
\chead{MTH 553: Lecture \# 19}
\rhead{Cliff Sun}

\begin{document}

\maketitle

\section*{Lecture Span}
\begin{enumerate}
    \item More heat equation
\end{enumerate}

\section*{Last time}
\begin{enumerate}
    \item Heat kernel $K(x,y,t)$ such that 
    \begin{equation}
        u(x,t) = \int_{\mathbb{R}^n} K(x,y,t) f(y) dy
    \end{equation}
    \item Maximum principle 
    \begin{equation}
        u(x,t) \leq \max_{(\tilde{x}, \tilde{t}) \in \Gamma} u(\tilde{x}, \tilde{t})
    \end{equation}
\end{enumerate}

\section*{Maximum Principle}

We now prove the Maximum principle for the heat equation.

\begin{proof}
    Assume that $u_t < \triangle u$ in $\Omega_T \cup \text{ top}$. At a maximum point in either $\Omega_T$ or the top $\Omega \times \{T\}$. Then at that maximum point, then $u_t \leq 0$ by the first
    derivative test and $\triangle u \leq 0$ by the 2nd derivative test. This contradicts $u_t < \triangle u$. Therefore, the maximum can only be on the bottom or sides. 
\end{proof}

Next, 

\begin{proof}
    Suppose that $u_t \leq \triangle u$ in $\Omega_T$. Then let $\epsilon > 0$, define $v = u - \epsilon t$ and $v_t = u_t - \epsilon \leq \triangle u - \epsilon = \triangle v  - \epsilon$. Therefore, 
    \begin{equation}
        v_t < \triangle v
    \end{equation}
    in $\Omega_T \cup \{\text{top}\}$. So therefore, 
    \begin{equation}
        v(x,t) \leq \max_{\Gamma}v \leq \max_{\Gamma}u
    \end{equation}
    Then $\epsilon \rightarrow 0$, then $v(x,t) \rightarrow u(x,t)$. 
\end{proof}

This proof was simply just abusing the inequality $u_t \leq u_{xx}$ and proving an inequality. 

\section*{Uniqueness forwards in time}

If $u,v \in C^{2,1}(\overline{\Omega_T})$, both then 

\begin{align*}
    w_t &= \triangle w \\
    w(x,0) &= g(x) \\
    w(x,t) &= h(x)
\end{align*}

Then $u \equiv v$. 

\begin{proof}
    Let $z = u - r$, so then $z_t = \triangle z$ on $\Omega_T$ and $z=0$ on $\Gamma$ by IC, BC. Therefore, 
    \begin{equation}
        z \leq 0 \quad \text{in} \Omega_T
    \end{equation}
    Similarly, for $-z$. Therefore, $z \equiv 0$. 
\end{proof}

\section*{Heat equation in $\mathbb{R}^n$}

\begin{align*}
    u_t &= \triangle u \\
    u(x,0) &= g(x)
\end{align*}

Tool, use the Fourier transform. Then, the Schwartz class 

\begin{equation}
    S = \{v \in C^{\infty}(\mathbb{R}^n): |x|^k|D^{\alpha}v(x)| \text{ in bounded on $\mathbb{R}^n$ for all $k \geq 0$}\}
\end{equation}

In particular, $C^{\infty}_{0} \subset S$. These are also functions that decay fast. Also, $e^{-|x|^2}$. Note, if $v \in S$, then 

\begin{enumerate}
    \item All derivatives of a Schwartz function are Schwartz functions. 
    \item $pv \in S$ if $p$ is a polynomial. 
    \item $|v(x)| \leq (\text{const.}/(1 + |x|)^{n+1})$, $v \in L^1(\mathbb{R}^n)$ and $v(x) \rightarrow 0$ as $x \rightarrow \infty$. 
\end{enumerate}

\begin{definition}
    For a integrable function $g \in L^1$, define its Fourier transform $\hat{g}$ by 
    \begin{equation}
        \hat{g}(\xi) = \frac{1}{\left( 2\pi \right)^{n/2}}\int_{\mathbb{R}^n}e^{-ix \cdot \xi} g(x) dx
    \end{equation}
    For $\xi \in \mathbb{R}^n$. 
\end{definition}

Then if $g(x) \in S$, then 

\begin{equation}
    \partial_{\xi_j}g(\xi) = \frac{1}{(2\pi)^{n/2}}\int_{\mathbb{R}^n}e^{-i x \cdot \xi}(-i x_j) g(x)dx
\end{equation}

This is justified by the "dominated convergence theorem". Then, $-ix_j g(x) \in L^1$. But we also showed that this is the same as 

\begin{equation}
    -i x_j \hat{g}(\xi)
\end{equation}

Repeating yields that 

\begin{equation}
    (D_{\xi})^{\beta}\hat{g}(\xi) = (-ix)^{\beta}\hat{g}(\xi)
\end{equation}

Note that 

\begin{equation}
    x^{\beta} = x^{\beta_1}x^{\beta_2}\dots
\end{equation}

Hence, $\hat{g} \in C_{0}^{\infty}(\mathbb{R}^n)$, or is smooth. Next, what if we do the following:

\begin{equation}
    \hat{\left( \frac{\partial g}{\partial x_i} \right)}(\xi) = \frac{1}{\left( 2\pi \right)^{n/2}}\int_{\mathbb{R}^n}e^{-ix \cdot \xi}\frac{\partial g}{\partial x_i}(x)dx
\end{equation}

THen we obtain 

\begin{equation}
    \frac{1}{(2\pi)^{n/2}}\int_{\mathbb{R}^n} = (i \xi_j)e^{-ix \cdot \xi}g(x)dx
\end{equation}

The boundary terms go away because $g$ is a Schwartz function. Therefore, we obtain 

\begin{equation}
    i\xi_j\hat{g}(x)
\end{equation}

Therefore, 

\begin{equation}
    D_x^{\alpha}\left( g \right)(\xi) = \left( -i\xi \right)^{\alpha}\hat{g}(x)
\end{equation}

So the Fourier transform assigns differentiation towards multiplication by polynomials. 

\begin{lemma}
    $$g \in L^1 \implies \hat{g} \in L^{\infty}$$
    That is $\hat{g}$ is bounded. Therefore, 
    $$g \in S \implies \hat{g} \in S$$

    That is $\hat{g}$ is smooth and rapidly decaying. 
\end{lemma}

\begin{proof}
    $g \in L^1$ and $$|\hat{g}(\xi)| \leq \int_{\mathbb{R}^n}\left( |g(x)| \right)dx < \infty$$ This is independent of $\xi$.
    This also means that $\hat{g} \in L^{\infty}$
\end{proof}

\begin{proof}
    Next, suppose $g \in S$, then for all $\alpha, \beta$, we have that 
    \begin{equation}
        D_{x}^{\alpha}\left( (-ix)^{\beta}g(x) \right) \in S \subset L^1
    \end{equation}
    Therefore, $\hat{g}$ is also a Schwartz function. 
\end{proof}

With

\begin{equation}
    D^{\alpha}_{X}\left( (ix)^{\beta}g(x) \right)^{\hat{}}(\xi) = (i\xi)^{\alpha}\left[ (-ix)^{\beta}g \right]^{\hat{}}(\xi)
\end{equation}
\begin{equation}
    = (i\xi)^{\alpha}D^{\beta}_{\xi}g(\xi)
\end{equation}

But the first statement is bounded, therefore, the last thing is also bounded. This is the condition that 

\begin{equation}
    |\xi^{\alpha}| |D^{\beta}_{\xi}g|
\end{equation}

is also bounded, therefore $g \in S$. 

\begin{definition}
    For $g \in L^1$, then define its inverse Fourier transform 
    \begin{equation}
        g^{\text{upside down hat}} = \frac{1}{(2\pi)^{n/2}}\int_{\mathbb{R}^n}e^{i x \cdot \xi}g(x)dx
    \end{equation}
    Note the plus sign change. This is the same as 
    \begin{equation}
        \hat{g}(-\xi)
    \end{equation}
\end{definition}

\end{document}